% Appendices. Reviewers need not read them; every claim must stand in the eight pages.
% Tables under tables/ are generated (gen_rule_tables.py, gen_prompt_appendix.py).
% Protocols follow docs/DESIGN_DECISIONS.md sections 10 and 11.

\section{Prompts}
\label{app:prompts}

Every prompt is rendered from a YAML template by the evaluation harness; the three paraphrases (p0, p1, p2) carry the same content in different wording and order, and the English-instruction ablation changes only the instruction language while the item and the answer marker \vi{đáp án:} stay Vietnamese. Demonstrations are built by the rule engine from three fixed base pairs (\vi{trung bình}, \vi{làm chủ}, \vi{vô hình}) whose syllables occur in no item; \task{T2} demonstrations are the same for every item because the prompt withholds the variant. No model receives a system prompt; the three demonstrations are inline in the single user turn. The examples below are the p0 templates with three demonstrations, exactly as rendered, for one example item per task; the sha256 of every prompt sent to a model is stored with its output. \placeholder{[NATIVE-CHECK] Every Vietnamese string awaits validation by a native speaker before the runs; the six-type overview and the region-neutral wording of DESIGN\_DECISIONS 3 are to be reflected in the templates.}

% GENERATED by paper/gen_prompt_appendix.py on 2026-09-30 17:05Z from prompts/noilai.yaml + prompts/demos.yaml; do not edit.
% prompt file sha256: demos.yaml=5fe710f99645, noilai.yaml=a390e6a44153, xcopa.yaml=de275fa2df4b, _all=002fb57ab3d6
\paragraph{Transformation (\task{T1}), variant \var{1}, paraphrase p0, three demonstrations.} Prompt id \texttt{vi-p0-s3-explained-raw}, sha256 \texttt{fc90f3c2182f}.
\begin{quote}\small
Nói lái là một cách chơi chữ trong tiếng Việt: hai âm tiết của một cụm từ trao đổi với nhau một số thành phần (phụ âm đầu, vần hoặc thanh điệu) để tạo thành một cụm từ mới. \\
Mỗi âm tiết tiếng Việt gồm ba phần: phụ âm đầu (có thể không có), vần (âm đệm nếu có, âm chính và âm cuối nếu có) và thanh điệu (một trong sáu thanh: ngang, huyền, sắc, hỏi, ngã, nặng). Ví dụ: \textquotedbl{}trung\textquotedbl{} gồm phụ âm đầu \textquotedbl{}tr\textquotedbl{}, vần \textquotedbl{}ung\textquotedbl{} và thanh ngang; \textquotedbl{}bình\textquotedbl{} gồm phụ âm đầu \textquotedbl{}b\textquotedbl{}, vần \textquotedbl{}inh\textquotedbl{} và thanh huyền.

Kiểu nói lái cần áp dụng: đổi vần, giữ phụ âm đầu và thanh điệu. \\
Cách làm từng bước: \\
1. Tách mỗi âm tiết thành phụ âm đầu, vần và thanh điệu. \\
2. Đổi vần của âm tiết thứ nhất cho vần của âm tiết thứ hai. \\
3. Giữ nguyên phụ âm đầu và thanh điệu của mỗi âm tiết ở vị trí cũ. \\
4. Ghép lại thành hai âm tiết và viết đúng chính tả. \\
Khi ghép lại, hãy viết đúng chính tả tiếng Việt: âm /k/ viết là \textquotedbl{}k\textquotedbl{} trước e, ê, i, viết là \textquotedbl{}q\textquotedbl{} trước âm đệm u và viết là \textquotedbl{}c\textquotedbl{} trong các trường hợp còn lại (ví dụ \textquotedbl{}c\textquotedbl{} ghép với vần \textquotedbl{}eo\textquotedbl{} phải viết là \textquotedbl{}keo\textquotedbl{}, không viết \textquotedbl{}ceo\textquotedbl{}); tương tự, \textquotedbl{}gh\textquotedbl{} và \textquotedbl{}ngh\textquotedbl{} đứng trước e, ê, i (ví dụ \textquotedbl{}nghe\textquotedbl{}, không viết \textquotedbl{}nge\textquotedbl{}); dấu thanh đặt đúng vị trí.

Ví dụ: \\
Cụm từ: trung bình \\
Đáp án: trinh bùng

Cụm từ: làm chủ \\
Đáp án: lù chảm

Cụm từ: vô hình \\
Đáp án: vinh hồ

Bây giờ hãy áp dụng kiểu nói lái trên cho cụm từ sau. \\
Cụm từ: mèo cái \\
Chỉ viết cụm từ kết quả trên một dòng riêng ở cuối, theo đúng dạng: \\
Đáp án: \textless{}cụm từ kết quả\textgreater{}
\end{quote}
\paragraph{Decoding (\task{T2}), paraphrase p0, three demonstrations (the variant is withheld).} Prompt id \texttt{vi-p0-s3-explained-raw}, sha256 \texttt{785246b4733d}.
\begin{quote}\small
Nói lái là một cách chơi chữ trong tiếng Việt: hai âm tiết của một cụm từ trao đổi với nhau một số thành phần (phụ âm đầu, vần hoặc thanh điệu) để tạo thành một cụm từ mới. Có bốn kiểu nói lái thường gặp: \\
- đổi vần, giữ phụ âm đầu và thanh điệu: Đổi vần của âm tiết thứ nhất cho vần của âm tiết thứ hai. Giữ nguyên phụ âm đầu và thanh điệu của mỗi âm tiết ở vị trí cũ. \\
- đổi cả phụ âm đầu và vần, giữ thanh điệu: Đổi chỗ phần phụ âm đầu và vần của hai âm tiết cho nhau (tức là đổi chỗ hai âm tiết). Giữ thanh điệu ở vị trí cũ: âm tiết đứng trước vẫn mang thanh của âm tiết đứng trước, âm tiết đứng sau vẫn mang thanh của âm tiết đứng sau. \\
- đổi thanh điệu, giữ phụ âm đầu và vần: Đổi thanh điệu của âm tiết thứ nhất cho thanh điệu của âm tiết thứ hai. Giữ nguyên phụ âm đầu và vần của mỗi âm tiết. \\
- đổi vần và thanh điệu, giữ phụ âm đầu: Đổi cả vần và thanh điệu của âm tiết thứ nhất cho vần và thanh điệu của âm tiết thứ hai. Giữ nguyên phụ âm đầu của mỗi âm tiết ở vị trí cũ. \\
Mỗi âm tiết tiếng Việt gồm ba phần: phụ âm đầu (có thể không có), vần (âm đệm nếu có, âm chính và âm cuối nếu có) và thanh điệu (một trong sáu thanh: ngang, huyền, sắc, hỏi, ngã, nặng). Ví dụ: \textquotedbl{}trung\textquotedbl{} gồm phụ âm đầu \textquotedbl{}tr\textquotedbl{}, vần \textquotedbl{}ung\textquotedbl{} và thanh ngang; \textquotedbl{}bình\textquotedbl{} gồm phụ âm đầu \textquotedbl{}b\textquotedbl{}, vần \textquotedbl{}inh\textquotedbl{} và thanh huyền.

Ví dụ: \\
Cụm từ: trinh bùng \\
Đáp án: trung bình

Cụm từ: chù lảm \\
Đáp án: làm chủ

Cụm từ: vồ hinh \\
Đáp án: vô hình

Cụm từ sau đây là một cách nói lái. Cụm từ gốc là gì? \\
Cụm từ: mài kéo \\
Chỉ viết cụm từ gốc trên một dòng riêng ở cuối, theo đúng dạng: \\
Đáp án: \textless{}cụm từ gốc\textgreater{}
\end{quote}
\paragraph{Validity (\task{T3}), variant \var{1}, paraphrase p0, a spelling twin as candidate.} Prompt id \texttt{vi-p0-s3-explained-raw}, sha256 \texttt{855d848b82af}.
\begin{quote}\small
Nói lái là một cách chơi chữ trong tiếng Việt: hai âm tiết của một cụm từ trao đổi với nhau một số thành phần (phụ âm đầu, vần hoặc thanh điệu) để tạo thành một cụm từ mới. \\
Mỗi âm tiết tiếng Việt gồm ba phần: phụ âm đầu (có thể không có), vần (âm đệm nếu có, âm chính và âm cuối nếu có) và thanh điệu (một trong sáu thanh: ngang, huyền, sắc, hỏi, ngã, nặng). Ví dụ: \textquotedbl{}trung\textquotedbl{} gồm phụ âm đầu \textquotedbl{}tr\textquotedbl{}, vần \textquotedbl{}ung\textquotedbl{} và thanh ngang; \textquotedbl{}bình\textquotedbl{} gồm phụ âm đầu \textquotedbl{}b\textquotedbl{}, vần \textquotedbl{}inh\textquotedbl{} và thanh huyền.

Kiểu nói lái đang xét: đổi vần, giữ phụ âm đầu và thanh điệu. \\
Cách làm từng bước: \\
1. Tách mỗi âm tiết thành phụ âm đầu, vần và thanh điệu. \\
2. Đổi vần của âm tiết thứ nhất cho vần của âm tiết thứ hai. \\
3. Giữ nguyên phụ âm đầu và thanh điệu của mỗi âm tiết ở vị trí cũ. \\
4. Ghép lại thành hai âm tiết và viết đúng chính tả. \\
Khi ghép lại, hãy viết đúng chính tả tiếng Việt: âm /k/ viết là \textquotedbl{}k\textquotedbl{} trước e, ê, i, viết là \textquotedbl{}q\textquotedbl{} trước âm đệm u và viết là \textquotedbl{}c\textquotedbl{} trong các trường hợp còn lại (ví dụ \textquotedbl{}c\textquotedbl{} ghép với vần \textquotedbl{}eo\textquotedbl{} phải viết là \textquotedbl{}keo\textquotedbl{}, không viết \textquotedbl{}ceo\textquotedbl{}); tương tự, \textquotedbl{}gh\textquotedbl{} và \textquotedbl{}ngh\textquotedbl{} đứng trước e, ê, i (ví dụ \textquotedbl{}nghe\textquotedbl{}, không viết \textquotedbl{}nge\textquotedbl{}); dấu thanh đặt đúng vị trí.

Ví dụ: \\
Cụm từ gốc: trung bình; cụm từ cần xét: trinh bùng \\
Đáp án: Có

Cụm từ gốc: làm chủ; cụm từ cần xét: chù lảm \\
Đáp án: Không

Cụm từ gốc: vô hình; cụm từ cần xét: vinh hồ \\
Đáp án: Có

Câu hỏi: Cụm từ \textquotedbl{}mài céo\textquotedbl{} có phải là kết quả đúng khi nói lái cụm từ \textquotedbl{}mèo cái\textquotedbl{} theo kiểu trên không? Kết quả đúng phải đúng cả phụ âm đầu, vần, thanh điệu và chính tả. \\
Trả lời \textquotedbl{}Có\textquotedbl{} hoặc \textquotedbl{}Không\textquotedbl{} trên một dòng riêng ở cuối, theo đúng dạng: \\
Đáp án: \textless{}Có hoặc Không\textgreater{}
\end{quote}


\section{Rule Tables}
\label{app:rules}

The tables in this appendix are generated from the parser's constants and the spelling lists, so that they cannot drift from the code: Table~\ref{tab:components} lists the canonical symbols, Table~\ref{tab:onsets} the surface spelling of every onset, Table~\ref{tab:spelling} the remaining spelling rules with engine-spelled examples, Table~\ref{tab:placement} the 69 syllables whose tone-mark letter differs between the two placement conventions, and Table~\ref{tab:variants} the six variants with engine-computed examples and the status of each (the four generated \task{T1} cells; \var{5}, the onset swap, and \var{6}, onsets and tones, are the order-reversed readings of \var{4} and \var{1} and enter the taxonomy, the \task{T2} gold and the \task{T3} twins).

% GENERATED by paper/gen_rule_tables.py (components) from noilai.vi.syllable constants and the Hunspell lists; generated 2026-09-30 15:54Z. Do not edit by hand.

\begin{table*}[t]\centering\small
\begin{tabular}{@{}lp{0.78\textwidth}@{}}
\toprule
Slot & Canonical symbols (as the parser writes them) \\
\midrule
Onset (24) & \texttt{$\varnothing$}, \texttt{b}, \texttt{c}, \texttt{ch}, \texttt{d}, \texttt{đ}, \texttt{g}, \texttt{gi}, \texttt{h}, \texttt{kh}, \texttt{l}, \texttt{m}, \texttt{n}, \texttt{ng}, \texttt{nh}, \texttt{p}, \texttt{ph}, \texttt{r}, \texttt{s}, \texttt{t}, \texttt{th}, \texttt{tr}, \texttt{v}, \texttt{x}. \texttt{c} stands for /k/ however spelled (c, k, q); \texttt{g} for the voiced velar fricative however spelled (g, gh); \texttt{ng} for the velar nasal however spelled (ng, ngh); \texttt{gi} is the onset that contracts a following i; \texttt{p} occurs only in loanwords. \\
Glide & present or absent (written \texttt{u} or \texttt{o}; the \texttt{u} of \texttt{qu} is the glide). \\
Nucleus (18) & \texttt{a}, \texttt{ă}, \texttt{â}, \texttt{e}, \texttt{ê}, \texttt{i}, \texttt{o}, \texttt{ô}, \texttt{ơ}, \texttt{u}, \texttt{ư}, \texttt{ia}, \texttt{iê}, \texttt{ua}, \texttt{uô}, \texttt{ưa}, \texttt{ươ}, \texttt{oo}. Open-only (never take a coda): \texttt{ia}, \texttt{ua}, \texttt{ưa}. Closed-only (always take a coda): \texttt{iê}, \texttt{uô}, \texttt{â}, \texttt{ă}, \texttt{ươ}. \\
Coda (11) & \texttt{$\varnothing$}, \texttt{m}, \texttt{n}, \texttt{ng}, \texttt{nh}, \texttt{p}, \texttt{t}, \texttt{c}, \texttt{ch}, \texttt{j}, \texttt{w}. \texttt{j} is the semivowel written i/y, \texttt{w} the semivowel written o/u; the stop codas \texttt{c}, \texttt{ch}, \texttt{p}, \texttt{t} take only the tones sắc and nặng. \\
Tone (6) & 0~=~ngang, 1~=~huyền, 2~=~sắc, 3~=~hỏi, 4~=~ngã, 5~=~nặng; index 0 (ngang) is the unmarked level tone, the other five carry a diacritic. \\
\bottomrule
\end{tabular}
\caption{The parser's canonical symbols (\texttt{noilai.vi.syllable}). A syllable is onset + rime + tone, the rime being glide$^?$ + nucleus + coda$^?$. Legality is decided against the attested inventory (Hunspell vi\_VN, both placement styles, extended with word-list syllables that reuse an attested rime).}
\label{tab:components}
\end{table*}

% GENERATED by paper/gen_rule_tables.py (onsets) from noilai.vi.syllable constants and the Hunspell lists; generated 2026-10-01 23:45Z. Do not edit by hand.

\begin{table*}[t]\centering\small
\begin{tabular}{@{}lll@{}}
\toprule
Onset & Written & Engine examples \\
\midrule
\texttt{$\varnothing$} & (nothing) & a, yên, ý \\
\texttt{b} & \texttt{b} & ba, bồng \\
\texttt{c} & \texttt{c, k, qu} & ca / ke / qua (c before a, k before e \^{e} i y, q before the glide) \\
\texttt{ch} & \texttt{ch} & cha, chồng \\
\texttt{d} & \texttt{d} & da, dang \\
\texttt{đ} & \texttt{đ} & đa, đồng \\
\texttt{g} & \texttt{g, gh} & ga / ghe (gh before e \^{e} i) \\
\texttt{gi} & \texttt{gi} & gia / gì / giêng (the rime's initial i is absorbed) \\
\texttt{h} & \texttt{h} & ha, hồng \\
\texttt{kh} & \texttt{kh} & kha, khang \\
\texttt{l} & \texttt{l} & la, lồng \\
\texttt{m} & \texttt{m} & ma, mồng \\
\texttt{n} & \texttt{n} & na, nồng \\
\texttt{ng} & \texttt{ng, ngh} & nga / nghe / ngoe (ngh before e \^{e} i; not with a glide) \\
\texttt{nh} & \texttt{nh} & nha, nhồng \\
\texttt{p} & \texttt{p} & pa, pin \\
\texttt{ph} & \texttt{ph} & pha, phồng \\
\texttt{r} & \texttt{r} & ra, rồng \\
\texttt{s} & \texttt{s} & sa, sồng \\
\texttt{t} & \texttt{t} & ta, tồng \\
\texttt{th} & \texttt{th} & tha, thang \\
\texttt{tr} & \texttt{tr} & tra, trồng \\
\texttt{v} & \texttt{v} & va, vồng \\
\texttt{x} & \texttt{x} & xa, xồng \\
\bottomrule
\end{tabular}
\caption{Surface spelling of the canonical onsets. The k/gh/ngh spellings are triggered by the first \emph{written} letter after the onset (e, \^{e}, i, y), so a glide switches them off (\texttt{ngoe}, \texttt{nguy}); every example is produced by \texttt{noilai.vi.syllable.spell}, and the examples of the plain rows are attested syllables of the release inventory.}
\label{tab:onsets}
\end{table*}

% GENERATED by paper/gen_rule_tables.py (spelling) from noilai.vi.syllable constants and the Hunspell lists; generated 2026-09-30 15:54Z. Do not edit by hand.

\begin{table*}[t]\centering\small
\begin{tabular}{@{}lp{0.42\textwidth}p{0.36\textwidth}@{}}
\toprule
Rule & Statement & Engine examples \\
\midrule
Glide letter & \texttt{u} after q and before \^{a} \^{e} \horn{o} \^{o} i ia i\^{e}; \texttt{o} before a \u{a} e & hoa, hoe, huê, tuân, qua \\
Nucleus i & written \texttt{y} after the glide and in the bare syllable; \texttt{i} elsewhere & quy, ý, lí, minh \\
Nucleus ia / i\^{e} & ya / y\^{e} after the glide; y\^{e} with a zero onset & quya, quyên, yêu, tiếng \\
Short \u{a} & written \texttt{a} before the semivowel codas & tay vs.\ tai; tau vs.\ tao \\
Coda /j/ & \texttt{y} after \u{a} \^{a}, \texttt{i} elsewhere & mây, mai, môi \\
Coda /w/ & \texttt{o} after a e, \texttt{u} elsewhere & chao, keo, mưu, chiêu \\
Stop codas & p t c ch take only sắc and nặng & mật, bất; *mât is rejected by the inventory \\
Open-only nuclei & ia ua \horn{u}a never take a coda & mía, mùa \\
Closed-only nuclei & i\^{e} u\^{o} \horn{u}\horn{o} \u{a} \^{a} always take a coda & muốn, mượn \\
Tone mark, new style & on the nucleus letter (second letter of i\^{e} y\^{e} u\^{o} \horn{u}\horn{o}, first of ia ua \horn{u}a) & tiếng, muốn, mía \\
Tone mark, old style & differs only on open oa oe uy after a non-q onset (mark on the first vowel letter) & hoà / hòa, khoẻ / khỏe, thuỷ / thủy \\
\bottomrule
\end{tabular}
\caption{Spelling rules the generator applies when it reassembles a syllable, with examples spelled by the engine. The same rules, run in reverse, are the parser's spelling round-trip: a syllable whose standard spelling does not reproduce the input is rejected in strict mode (so \vi{c\'{e}o} for \vi{k\'{e}o} scores as a spelling error, never as correct).}
\label{tab:spelling}
\end{table*}

% GENERATED by paper/gen_rule_tables.py (placement) from noilai.vi.syllable constants and the Hunspell lists; generated 2026-09-30 17:22Z. Do not edit by hand.

% 69 pairs (new style, old style)
\begin{table}[t]\centering\small
\begin{tabular}{@{}llllll@{}}
\toprule
New & Old & New & Old & New & Old \\
\midrule
choá & chóa & luỹ & lũy & toè & tòe \\
choé & chóe & ngoã & ngõa & toé & tóe \\
choẹ & chọe & ngoé & ngóe & toạ & tọa \\
chuỳ & chùy & nguỵ & ngụy & toả & tỏa \\
doá & dóa & nhoà & nhòa & toẻ & tỏe \\
doạ & dọa & nhoá & nhóa & toẽ & tõe \\
goá & góa & nhoè & nhòe & truỵ & trụy \\
hoà & hòa & nhoé & nhóe & tuý & túy \\
hoá & hóa & nhuỵ & nhụy & tuỳ & tùy \\
hoè & hòe & oà & òa & tuỵ & tụy \\
hoạ & họa & oé & óe & tuỷ & tủy \\
hoả & hỏa & oẹ & ọe & uý & úy \\
hoẹ & họe & oẻ & ỏe & uỷ & ủy \\
huý & húy & soà & sòa & xoà & xòa \\
huỷ & hủy & suý & súy & xoá & xóa \\
khoá & khóa & thoà & thòa & xoã & xõa \\
khoé & khóe & thoá & thóa & xoè & xòe \\
khoả & khỏa & thoả & thỏa & xoả & xỏa \\
khoẻ & khỏe & thuý & thúy & xoẹ & xọe \\
loà & lòa & thuỳ & thùy & xuý & xúy \\
loè & lòe & thuỵ & thụy & xuỳ & xùy \\
loé & lóe & thuỷ & thủy & đoá & đóa \\
luỵ & lụy & toà & tòa & đoạ & đọa \\
\bottomrule
\end{tabular}
\caption{The 69 syllables of the Hunspell vi\_VN list whose tone-mark letter differs between the new placement convention (mark on the nucleus letter) and the old one (mark on the first vowel letter of the open rimes oa, oe, uy). The set is the difference between the two released lists, and the engine's rule reproduces every pair; the C2 arm converts running text between the two conventions and changes exactly these syllables (and their capitalized forms).}
\label{tab:placement}
\end{table}

% GENERATED by paper/gen_rule_tables.py (variants) from noilai.vi.syllable constants and the Hunspell lists; generated 2026-10-01 23:45Z. Do not edit by hand.

\begin{table*}[t]\centering\small
\begin{tabular}{@{}lllll@{}}
\toprule
Variant & Swapped & Kept in place & Engine example & Status \\
\midrule
\var{1} & rime & onset, tone & \vi{mèo cái} $\rightarrow$ \vi{mài kéo} & generated cell \\
\var{2} & onset, rime & tone & \vi{đầu tiên} $\rightarrow$ \vi{tiền đâu} & generated cell \\
\var{3} & tone & onset, rime & \vi{bí mật} $\rightarrow$ \vi{bị mất} & generated cell \\
\var{4} & rime, tone & onset & \vi{bí mật} $\rightarrow$ \vi{bật mí} & generated cell \\
\var{5} & onset & rime, tone & \vi{giải pháp} $\rightarrow$ \vi{phải giáp} & taxonomy; \task{T2} gold; \task{T3} twins \\
\var{6} & onset, tone & rime & \vi{mèo cái} $\rightarrow$ \vi{kéo mài} & taxonomy; \task{T2} gold; \task{T3} twins \\
\bottomrule
\end{tabular}
\caption{The six \vi{nói lái} variants as component permutations of a syllable pair (\texttt{noilai.gen.variants}): the non-empty proper subsets of \{onset, rime, tone\} exchanged between the two positions. Each is an involution; reading an output in the other order maps \var{1} to \var{6}, \var{2} to \var{3} and \var{4} to \var{5}, so \var{5} and \var{6} add no new unordered pair and are not generated as \task{T1} cells. Examples \var{1}--\var{5} are attested folk forms that the engine reproduces exactly; the \var{6} example is the \var{1} pair read in the other order, and the seed's two attested \var{6} rows (\vi{thay đổi} $\rightarrow$ \vi{đảy thôi}, \vi{trái gió} $\rightarrow$ \vi{giái tró}) are reproduced exactly as well.}
\label{tab:variants}
\end{table*}


\section{Validation Protocol and Human Baseline}
\label{app:validation}

\paragraph{Validation.}
The packet has five parts, judged in order of priority. (A) Calibration: 16 items keyed by the rule engine on pairs that are not in the release (four plain correct \task{T1} items, five planted errors, and the cases the instructions single out: new-style placement, \vi{y} after a consonant, a vulgar output, a \vi{hỏi}/\vi{ngã} output, a \task{T2} item, a correct \task{T3} item and a \task{T3} spelling twin); a validator who matches the key on fewer than 80\% of the correctness answers is briefed again and takes a second calibration set; nobody is excluded on calibration. (B) Generated items: 30 per task $\times$ variant cell drawn as a stratified probability sample (strata: cell $\times$ lexical or pseudo source, proportional allocation by largest remainder; the population size, sample size and weight of every stratum are recorded; \task{T3} items are drawn singly, never both members of a twin pair) plus four planted controls per cell built from items outside the sample: a \task{T1} candidate replaced by another variant's output (never the gold in the other order, which is correct leniently) or by the gold with one tone changed, a \task{T2} candidate that is no reading of the input under any of the six variants in either order, and a \task{T3} item shown with the opposite verdict. The 408 rows are shuffled under opaque identifiers; 60 go to every validator and the rest to two validators in rotating pairs (all rows to both when there are two). Validators are told that some rows are deliberately wrong. (C) Every attested row, seed and web-sourced candidates alike, judged by every validator for validity, familiarity, spelling and offensiveness; a row counts as native-verified when at least two validators accept it and none rejects it, and it enters the attested set only with its citation checked on the source page. (D) Twenty pairs on which the validator writes the nói lái they produce first (onsets, rimes and tones all differ, so the answer identifies the variant), 20 \vi{qu}- pairs shown under both analyses, and 14 spelling-convention questions (tone-mark placement, \vi{i/y}, \vi{oao}, \vi{gi}). (E) Each core \task{T2} item's dictionary readings to accept or reject, with any missing reading added; 50 items are judged by every validator. For each Part B row the validator answers with \emph{yes}, \emph{no} or \emph{unsure}: (1) \emph{correct}: the output is the correct nói lái of the input under the named variant (for \task{T2}, the input is a nói lái of the candidate original under some variant and order; for \task{T3}, the system's verdict is right); (2) \emph{spelling}: the output is spelled acceptably, either tone-mark placement and the \vi{lí/lý} alternation being acceptable; (3) \emph{lexical}: the input, and separately the output, is a real Vietnamese word or phrase (the speaker's judgment, not dictionary membership); (4) \emph{offensive}: the input or the output has an offensive or vulgar reading, including one reached by decoding; and (5) \emph{dialect}: the regional merger, if any, on which the correctness answer depends. Each validator's region of upbringing is recorded. Agreement is Krippendorff's nominal $\alpha$ on yes/no with ``unsure'' as missing (as a third category in a sensitivity analysis), Gwet's AC1, raw pairwise agreement and the label marginals, each with a 1,000-replicate item bootstrap interval; $\alpha$ on correctness is computed over the Part B sheet including its planted control rows, and the probability-sample-only $\alpha$ and AC1 are reported beside it; Part E uses MASI and Jaccard distances. Adjudication: unanimous definite labels stand; two validators who disagree are resolved by the third, who judges blind, by strict majority; a remaining split, or a disagreement with only two validators, is decided by the authors against the rule tables with a logged reason; an item with fewer than two definite labels is unresolved, stays in the benchmark and is left out of the precision estimate; the authors never override validators who agree. One validator's offensive flag is enough to remove an item from every API prompt and the human forms. A validator who labels fewer than 75\% of the planted controls ``no'' is briefed again after the first weekly return, and the rate is reported; no validator's labels are excluded or down-weighted on it. A comment that points at a rule is checked against the rule tables, and a confirmed generator bug triggers a fix and a full regeneration before any model run, which is reported. Generator precision per cell is $\sum_h W_h \hat p_h$ over the cell's strata with population shares $W_h$, with an unweighted Wilson interval beside it; the pooled estimate weights cells by population. The instructions given to validators, in Vietnamese and English, the consent form and the adjudication log are released with the data. \placeholder{Validators: number, regions, hours; per-judgment agreement table; adjudication counts; generator precision per cell.}

\paragraph{Human baseline.}
About twenty adult native speakers who did not take part in the validation each complete one form of 30 items (30--45 minutes): six anchor items seen by everyone (one per generated variant, one \task{T2}, one \task{T3}) and 24 items assigned so that every non-anchor item is seen by exactly two people, 246 items in all, about twenty per cell, drawn from the open-model main sample so that every model is scored on exactly these items; items flagged vulgar are excluded. Each row shows the exact p0 user message the models receive for that item (instruction, the same three demonstrations, the item and the answer marker) with its prompt hash recorded, in randomized item order, with one instruction-check row and a closing 10-item block of \task{T2} items on attested, non-vulgar examples that serves as the human reference for the cultural task; respondents write the answer (\task{T1}, \task{T2}) or tick yes/no (\task{T3}). They are asked not to use a dictionary, a search engine or an AI assistant and to say whether they did, and whether they validated items for this study. Answers are scored by the same code as the models, strictly, leniently and tolerant of doublets and homophones. Item-level human accuracy is the mean of its two judgments; the comparator is mean-human accuracy, with ``any human correct'' as a ceiling reported separately and an interval from a two-way bootstrap over persons and items; results are reported by output lexicality and by region, and the agreement between the two raters of an item is reported as nominal $\alpha$ on the produced answer and on correctness. A form is excluded, and counted, when fewer than half of its items are answered, when the instruction check fails, or when the respondent reports using a tool or having validated. If more volunteers materialize, items receive a third rater before the item set widens.

\section{Data Statement}
\label{app:datastatement}

Following \citet{Q18-1041}. \emph{Curation rationale}: items are generated by a rule engine from two-syllable entries of a Vietnamese word list and from sampled attested syllables, to test whether models can manipulate the components of a syllable; attested examples are collected from published collections of folk and literary nói lái, each with its source. \emph{Language variety}: standard written Vietnamese (\texttt{vi-VN}), lowercase, NFC, stored in the majority placement convention, \vi{i/y} per the stated rule; the attested subset includes Southern forms. \emph{Speaker demographics}: not applicable to generated items; attested items are folk material of unknown authorship or literary quotations with their source cited; the attested curation reflects the author's region. \emph{Annotator demographics}: three adult native speakers and about twenty human-baseline respondents, described in aggregate by age band, region grown up in, years in Vietnam and whether raised abroad (\placeholder{reported at submission}). \emph{Speech situation}: written wordplay, out of context. \emph{Text characteristics}: two-syllable phrases (three in a few attested cases), no punctuation, with per-syllable structure, lexicality and frequency covariates. \emph{Provenance}: the spelling list and word list (GPLv2), XCOPA (CC~BY~4.0), per-row citations for attested items, the generator commit and content hash (the build seed is withheld), resource hashes; the SHA-256 of the canary GUID \placeholder{printed here at release}, the GUID itself living in the gated README, the dataset card and the file header; split sizes and vulgar-flag counts; the rule-versus-folk mismatches and out-of-design rows of the attested set; the list entries the parser rejects, listed and explained; known limitations (\vi{qu}, regional labels, the \vi{i/y} convention); and the full parse table of the syllable inventory as supplementary material. The long-form statement is in the repository.

\section{Pre-registration}
\label{app:prereg}

The pre-registration is two-stage. The generator, data build, sampling files, scoring, extraction, census, placement baseline and constants are committed first (commit \placeholder{stage-1 hash}); the hypotheses of Sections~\ref{sec:results}--\ref{sec:tone}, the metrics, exclusion rules, stopping rules, analysis plan and probing protocol are fixed before any test-split run (the pilot runs on development items only) and committed on \placeholder{date of the stage-2 commit} (commit \placeholder{stage-2 hash}); deviations are logged with a date in the repository. Table~\ref{tab:hypotheses} restates them with their directional predictions and the quantity that decides each.

\begin{table*}[t]
  \centering\small
  \begin{tabular}{@{}lp{0.44\textwidth}p{0.38\textwidth}@{}}
    \toprule
    & Prediction & Decided by \\
    \midrule
    \hyp{1} & Alignment carries information about \task{T1} accuracy beyond count: $\beta(\text{align}) > 0$, $\beta(\text{tokens}) < 0$, and adding the split indicator and alignment to a model with tokens per syllable improves fit. & Nested likelihood-ratio test / $\Delta$AIC (model fixed effects, base-pair cluster-bootstrap CIs for $\beta$); a model with fewer than 300 misaligned split syllables contributes descriptively only, and fewer than four families clearing the floor demotes \hyp{1} to descriptive. \\
    \hyp{2} & Descriptive case study, not a tested hypothesis: PhoGPT-4B-Chat's whole-syllable vocabulary is the panel's extreme on the share of whole-syllable tokens, an E2 covariate; its per-syllable NLL and \task{T1} accuracy are reported beside Qwen3.5-4B and Gemma~3 4B. & Reported with CIs; the whole-syllable-share coefficient in E2 is the population-level test; no inequality is pre-registered. \\
    \hyp{3} & In pass-through tokenizers, $\tau(\mathrm{NFD})$ on \task{T1} \var{3} $> 0$ and on XCOPA $< 0$; distribution shift is the default explanation the crossover must beat. & Two paired tests per model, Holm within the pair; both family-level pooled CIs on the predicted side, the per-family table primary; below five informative families the pooled CI is labelled unreliable and the per-family CIs decide; a uniform drop is reported as OOD sensitivity. \\
    \hyp{3b} & The NFD gain on \var{3} is larger where the tone mark is isolated than where it stays attached (DiD $> 0$); the \var{1} contrast $\approx 0$. & Matched items, cluster bootstrap by base pair (McNemar mid-$p$ as the secondary item-level statistic), Holm over the four contrasts. \\
    \hyp{4} & The rarer placement convention (expected: new style, fixed by a corpus count before any run) lowers accuracy on affected items in models generally. & Per-model paired test on affected items, Holm across the panel; $k$ of $N$ plus the pooled estimate; TOST with $\delta = 2$ points for nulls. \\
    \hyp{5} & In a Gemma~3 model with \var{3} accuracy $< 25\%$, tone is decodable (selectivity and excess) at the token after the syllable by layer $L/2$, and, on the pairs whose clean answer is correct, the perception readout recovers under patching where the manipulation readout does not. & Section~\ref{sec:tone} protocol; $\delta_{\mathrm{sel}} = 0.15$ provisional; syllable-clustered CIs; below 50 retained pairs, the perception readout alone and a localization study; no earlier-under-NFD claim where the structural baseline is near 100\%. \\
    \hyp{6} & Exact attested items outscore matched generated items (pooled); guided-instruction contamination $> 0$ on attested, $\approx 0$ on matched generated. & Matching on variant, task, lexicality, trigger and tone class; per-model tests Holm across the panel plus pooled; the contamination discriminators carry the inference; floor: at least 100 verified exact rows from at least three collections, else descriptive. \\
    \bottomrule
  \end{tabular}
  \caption{Pre-registered hypotheses, their directional predictions and the quantity that decides each; ``pooled'' is a hierarchical random-effects estimate with models nested in tokenizer families, reported with the number of families, never a count of significant models.}
  \label{tab:hypotheses}
\end{table*}

\section{Compute, Environment and Panel Details}
\label{app:compute}

\placeholder{Per-run table: model, checkpoint id and revision, engine and version, requested and resolved precision, quantization, hardware, wall time, GPU-hours; totals for checklist item C1. Tokenizer profile and census verdict of every panel model (tokens per syllable, single-token share, alignment among split syllables, onset$|$rime split, NFD tone isolation, passes through / normalizes / corrupts per engine). Package versions (C4). Models that failed smoke tests and why. The adult account holder's role for the API runs.}

\section{Additional Results}
\label{app:extra}

\placeholder{Accuracy range over paraphrases per model; lenient \task{T1}; the spelling-twin column of \task{T3}; bf16 drift on about 200 items per quantized model; the two-engine C1 agreement on Gemma~3 1B; the reasoning-cost curve (accuracy against reasoning tokens on 500 items); per-variant results restricted to non-degenerate items; per-stratum tables; the realized design effect and per-model ICC; the \vi{qu}- stratum if its convention is settled in time.}
